% Pure Math Survey 2.1.0: paper. Preparation is not completed authorship.
% r5 opening: concise research abstract and Introduction; later count is editable.
% Current r5 authoring duties: active 2--3 normal TYPESET-line target for explanatory prose.
% Inspect the actual final PDF; record specific long-block reasons, alternatives,
% page/source positions, final PDF hash and reinspection. No TeX line-count games.
% Reconstruct obligations before locating listed steps; save actual fast-route,
% slow-step and condition-change attempts. Exact deep imports still need their
% current hypothesis checks, substitutions and object conversions explained.
% Reusable explanations return to live KB, review and a new snapshot before adoption.
% This scaffold is unfinished; profile/view tasks and substantive evidence remain.
\documentclass[12pt,reqno]{amsart}
\usepackage{math-review}
\input{../generated/conventions.tex}
\title{@@TITLE@@}
\date{}
\begin{document}
\maketitle
@@DRAFT@@
\begin{abstract}
\placeholder{Briefly state the problem, decisive hypotheses, actual contribution and key mechanism honestly. Keep literature comparison and proof detail for Introduction. Unauthorized research inputs, unverified contributions or incomplete proofs remain DRAFT_UNVERIFIED; no candidate result or conjecture is described as proved.}
\end{abstract}
\section{Introduction}
\placeholder{Read the selected reviewed overview, routes and node fast passages first; form Question, Main Claim, Mechanisms and Critical Conditions. Use full fixed canonical inputs below. For each borrowed principal place controlling source, version and exact locator in adjacent visible prose; distinguish adapted, local and unverified claims. Paper prior-work and novelty assertions need actual source and slow-detail verification.}
@@PRINCIPALS@@
\placeholder{Identify closest prior work and the specific unresolved difficulty; state complete results and precise supported differences, then explain the new idea and proof route. Match originality claims to authorized plan contributions and verified prior results; attribute known results honestly. Do not merely repeat the abstract. Later sections may split or combine to fit needed definitions, exact imports, new constructions and rigorous proofs; no fixed count.}
\section{Setup and Imported Inputs}
\placeholder{Open only the precise slow subsections required by actual proof promises. Establish consumed important Lemma/Proposition interfaces before their first proof use, separating independent branches. Keep all statement types numbered with labels and no descriptive optional titles; proof-target headings are allowed. Attribute borrowed, modified and composite tools nearby, with actual conversions and every source input; local derivations need no invented citation. Missing mathematics returns to KB, review and a new snapshot.}
@@INGREDIENTS@@
\placeholder{Only needed definitions and accurately attributed external theorems; establish reader bridges before first use and explain imported results' present hypothesis checks, substitution and output use.}
\section{Main Construction}
@@VIEW_TASK@@
\placeholder{Present the substantive new estimates or lemmas, with all quantifiers and constant dependencies.}
\section{Proofs of the Main Results}
\placeholder{Independently reconstruct the obligations required by the exact claims before checking the supplied proof list. Write complete proofs of all original principals. Disclose every unresolved gap; an incomplete contribution remains an internal draft. Save actual fast-route, slow-step and condition-change checks. Template completion or successful compilation cannot verify genuine originality or complete research-proof quality.}
\placeholder{Prefer 2--3 normal typeset lines for ordinary explanatory prose; keep continuous arguments together when needed. Rewrite hidden mathematical actions; record page/source positions, specific reasons and alternatives for necessary long blocks, the final PDF hash and reinspection scope. Never use decorative formulas, forced breaks, smaller type or omitted hypotheses.}
\bibliographystyle{plain}
\bibliography{../generated/references}
\end{document}
